\documentclass[12pt,a4paper]{ctexart}

% Packages
\usepackage{geometry}
\geometry{left=2.5cm,right=2.5cm,top=2.5cm,bottom=2.5cm}
\usepackage{graphicx}
\usepackage{booktabs}
\usepackage{longtable}
\usepackage{array}
\usepackage{multirow}
\usepackage{enumitem}
\usepackage{listings}
\usepackage{xcolor}
\usepackage{amsmath}
\usepackage{amssymb}
\usepackage{hyperref}
\usepackage{float}
\usepackage{caption}
\usepackage{subcaption}
\usepackage{fancyhdr}
\usepackage{algorithm}
\usepackage{algpseudocode}

% Code listings style
\lstset{
    language=Python,
    basicstyle=\ttfamily\small,
    keywordstyle=\color{blue}\bfseries,
    commentstyle=\color{gray},
    stringstyle=\color{orange},
    numbers=left,
    numberstyle=\tiny\color{gray},
    frame=single,
    breaklines=true,
    showstringspaces=false,
    tabsize=4,
    captionpos=b,
    escapeinside={(*@}{@*)}
}

% Hyperref setup
\hypersetup{
    colorlinks=true,
    linkcolor=blue,
    filecolor=magenta,
    urlcolor=cyan,
    citecolor=green,
}

% Header and footer
\pagestyle{fancy}
\fancyhf{}
\fancyhead[L]{\leftmark}
\fancyhead[R]{\thepage}
\renewcommand{\headrulewidth}{0.4pt}

% Title info
\title{\textbf{基于CMA设计思想的Meta-Harness Engineering框架改进研究}\\\vspace{0.5em}
\large A Study on MHE Framework Improvements Inspired by Claude Managed Agents Architecture}
\author{Meta-Harness Engineering Research Group}
\date{\today}

\begin{document}

\maketitle

\begin{abstract}
随着大语言模型（LLM）能力的迅速提升，基于LLM的智能体（Agent）系统已成为人工智能应用落地的核心形态。Meta-Harness Engineering（MHE）框架作为一套面向工程落地的元Harness智能体自我重长运行时，已经在组件化SDK、图模型验证、热加载和安全控制等方面建立了坚实的基础。与此同时，Anthropic于2026年推出并进入 public beta 的 Claude Managed Agents（CMA）代表了托管Agent基础设施的重要实践。根据其公开文档，CMA提供了托管Agent运行时、会话管理、权限策略与凭证保险库等能力；从工程分析角度，可将其理解为一种强调控制面与执行面解耦的设计范式。

本文系统分析CMA的核心架构原则，深入评估MHE框架的当前实现状态，识别出状态持久化、Harness无状态化、沙箱隔离、凭证管理和故障恢复等关键差距，并在此基础上提出七项具体的架构改进方案及分阶段实施路线图。研究表明，MHE框架通过引入CMA的设计思想，可以从"功能丰富的单体运行时"演进为"云原生的弹性基础设施"，从而更好地支撑生产级Agent系统的部署与运维。

\textbf{关键词：} Meta-Harness Engineering；Claude Managed Agents；Agent基础设施；云原生架构；自我改进智能体
\end{abstract}

\tableofcontents
\newpage

%==============================================================================
\section{引言}\label{sec:intro}
%==============================================================================

\subsection{研究背景}

随着大语言模型（Large Language Models, LLMs）能力的迅速提升，基于LLM的智能体（Agent）系统已成为人工智能应用落地的核心形态。从简单的单轮对话到复杂的多步骤任务调度、工具调用与多智能体协作，Agent系统的结构正在变得越来越复杂。然而，当前大多数Agent框架仍采用"固定配置"的设计范式：开发者根据经验手工编排提示词（Prompt）、工具链、记忆模块与控制流，一旦部署完成，其结构在运行期间基本保持不变。这种范式在面对动态变化的任务分布、性能瓶颈或资源约束时，往往显得缺乏弹性。

近年来，学术界和工业界开始关注"元层面"的优化——即不再仅仅优化模型权重，而是优化包裹在模型外部的运行框架（Harness）。斯坦福与MIT联合提出的Meta-Harness框架表明：在固定底层模型的情况下，仅通过优化harness代码，就可能在同一基准上产生高达6倍的性能差距。与此同时，英属哥伦比亚大学提出的ADAS（Automated Design of Agentic Systems）则展示了利用元智能体自动编写和迭代优化Agent代码的可行性。这些研究共同指向一个趋势：\textbf{Agent系统的设计本身应当成为可自动优化的对象}。

在这一背景下，Anthropic于2026年4月发布并进入 public beta 的 Claude Managed Agents（CMA）标志着托管 Agent 基础设施进入新的发展阶段。CMA并非一个新的模型，也不是Claude Code CLI的功能扩展，而是一个独立的API层面产品——开发者通过公开文档中给出的Agent与Session接口与Anthropic托管运行时交互，由平台负责会话管理、权限控制以及部分运行时基础设施。就产品直觉而言，CMA的核心价值主张可以概括为：\textbf{开发者定义Agent的任务、工具与护栏，平台承担更大比例的运行时管理职责}。

\subsection{研究目标与意义}

本研究的目标是将CMA的先进架构设计思想引入MHE框架的工程实现中，具体包括：

\begin{enumerate}
    \item \textbf{架构解耦}：借鉴CMA的Session/Harness/Sandbox三层解耦模型，将MHE从单体运行时重构为分层基础设施；
    \item \textbf{状态外部化}：实现持久化的会话事件日志，使Harness实例无状态化，支持水平扩展和故障恢复；
    \item \textbf{安全增强}：引入凭证保险库（Credential Vault）和架构级安全隔离，超越当前的最小化身份边界；
    \item \textbf{可观测性升级}：建立基于结构化事件流的统一可观测性模型，满足审计和合规需求。
\end{enumerate}

本研究的意义在于：为MHE框架从"研究和工程运行时"向"生产级控制平面"演进提供明确的架构路线图，同时验证"大模型厂商完全有能力将Harness作为托管服务提供，形成模型+Harness一体化竞争优势"这一行业趋势。

\subsection{研究方法与资料边界}

本文采用“代码与文档实证 + 托管Agent产品对照分析”的方法：对MHE部分，结论尽量直接建立在本仓库中的\texttt{README}、\texttt{USER\_GUIDE}、\texttt{TECHNICAL\_MANUAL}及核心源码之上；对CMA部分，则优先依据Anthropic公开文档、定价页、权限策略页与工程文章进行归纳。凡是能够被公开资料直接支持的内容，本文尽量按“事实陈述”表述；凡是为了与MHE进行结构对比而引入的抽象（如分层视角、控制面/执行面划分），本文均将其视为\textbf{工程分析模型}，而非对Anthropic内部实现细节的断言。

\subsection{论文结构}

本文共分为七章。第\ref{sec:mhe}节介绍MHE框架的核心概念与当前实现；第\ref{sec:cma}节深入解析CMA的架构设计原则；第\ref{sec:compare}节进行MHE与CMA的对比分析；第\ref{sec:improvements}节提出具体的改进方案；第\ref{sec:roadmap}节给出实施路线图；第\ref{sec:conclusion}节总结全文并展望未来研究方向。

%==============================================================================
\section{MHE框架概述}\label{sec:mhe}
%==============================================================================

\subsection{设计哲学}

Meta-Harness Engineering（MHE）是一套面向工程落地的元Harness智能体自我重长框架。其设计哲学深受数值偏微分方程（PDE）求解器中的"迭代收敛"思想影响，将Agent系统抽象为由九大核心基础组件加一个元层优化器（Optimizer）构成的可配置图结构，并通过性能瓶颈驱动的自我重长机制动态调整组件连接参数乃至生成新组件。

MHE框架特别强调\textbf{接口标准化}、\textbf{安全可控}与\textbf{工程可落地性}，试图为通用智能体的设计提供一套兼具灵活性与稳定性的新范式。框架的核心设计原则包括：

\begin{itemize}
    \item \textbf{声明式组件模型}：每个组件由JSON清单（Manifest）和Python实现共同定义，通过HarnessAPI收集声明；
    \item \textbf{候选图优先}：运行时从不直接修改活动图，而是通过暂存（stage）→验证（validate）→提交（commit）的流程确保变更安全；
    \item \textbf{内部图权威}：XML仅作为导入/配置格式，内部图快照（GraphSnapshot）才是运行时真理；
    \item \textbf{显式边界}：热加载、身份、安全和安全边界都是显式的运行时服务，而非隐藏的框架魔法。
\end{itemize}

\subsection{核心架构组件}

MHE的当前实现包含以下核心子系统，如表~\ref{tab:mhe-components}所示。

\begin{table}[H]
\centering
\caption{MHE核心子系统一览}\label{tab:mhe-components}
\begin{tabular}{@{}lll@{}}
\toprule
\textbf{子系统} & \textbf{核心职责} & \textbf{关键模块} \\
\midrule
SDK & 类型化组件接口、清单模型、运行时注入 & \texttt{sdk/} \\
发现与加载 & 清单解析、静态验证、依赖排序、实例化 & \texttt{sdk/discovery.py} \\
注册表 & 声明跟踪、插槽索引、能力索引、图指针 & \texttt{sdk/registry.py} \\
图核心 & 节点、边、快照、候选集、验证报告 & \texttt{core/models.py} \\
连接引擎 & 候选暂存、提交、回滚、路由表编译 & \texttt{core/connection\_engine.py} \\
引导编排器 & 组合发现、激活、处理器注册 & \texttt{core/boot.py} \\
安全链 & 沙箱验证、A/B影子测试、策略否决、自动回滚 & \texttt{safety/} \\
热加载 & 检查点、迁移适配器、交换编排、观察窗口 & \texttt{hotreload/} \\
身份边界 & 受保护的证明和凭证分离 & \texttt{identity/} \\
可观测性 & 指标、追踪、轨迹存储、PROV证据 & \texttt{observability/} \\
审计 & 仅追加审计日志、Merkle锚定 & \texttt{provenance/} \\
优化器 & 触发器、适应度、收敛、搜索阶段 & \texttt{optimizer/} \\
\bottomrule
\end{tabular}
\end{table}

\subsection{引导与运行时流程}

MHE的引导编排器\texttt{HarnessRuntime}执行以下主要流程：

\begin{enumerate}
    \item 从配置的 Discovery Root 解析清单；
    \item 应用启用/禁用覆盖；
    \item 运行静态清单验证；
    \item 解析依赖顺序；
    \item 实例化每个组件并通过\texttt{HarnessAPI}收集声明；
    \item 将声明注册到注册表；
    \item 收集声明的迁移适配器到共享运行时注册表；
    \item 使用\texttt{ComponentRuntime}激活组件；
    \item 向\texttt{ConnectionEngine}注册声明的连接处理器。
\end{enumerate}

值得注意的是，\textbf{引导不提交图}。图提交是一个独立的显式操作，通过\texttt{commit\_graph()}方法完成。

\subsection{当前成熟度与局限性}

MHE框架目前在以下方面表现成熟：

\begin{itemize}
    \item 组件SDK的类型安全性和声明式API；
    \item 图模型的候选优先变更管理和版本控制；
    \item 热加载协议的状态迁移适配器机制；
    \item 安全链的四层防御结构。
\end{itemize}

然而，当前实现仍存在以下关键局限，这些局限正是本研究试图解决的问题：

\begin{enumerate}
    \item \textbf{内存中的状态存储}：大多数有状态服务（注册表、版本存储、事件总线）均为内存实现，进程崩溃导致状态丢失；
    \item \textbf{有状态的Harness运行时}：\texttt{HarnessRuntime}持有活动组件实例字典，无法在不同进程间共享会话；
    \item \textbf{沙箱抽象强于实现}：沙箱层级结构存在，但默认回退到进程内适配器，未形成真正的安全边界；
    \item \textbf{身份边界最小化}：\texttt{InMemoryIdentityBoundary}仅提供基本的证明签发和凭证隔离，缺乏企业级的凭证保险库功能；
    \item \textbf{缺乏自动故障恢复}：热加载支持回滚，但无自动检查点和\texttt{wake(session\_id)}式的故障恢复机制；
    \item \textbf{可观测性有限}：\texttt{trace\_id}主要作为关联字段，未驱动采样或路由行为，事件流不完整。
\end{enumerate}

%==============================================================================
\section{CMA架构解析}\label{sec:cma}
%==============================================================================

\subsection{产品定位}

Claude Managed Agents（简称CMA）是Anthropic在Claude Platform上推出的一套可组合API服务，旨在让企业开发者无需自建agent loop、沙盒执行层或状态管理系统，即可大规模部署云托管的AI智能体。

CMA并非一个新模型，也不是Claude Code CLI的功能扩展，而是一个独立的API层面产品。根据公开文档，开发者可以通过创建Agent与Session的接口与Anthropic托管运行时交互；平台负责会话管理、权限策略以及部分运行时基础设施。

CMA的核心价值主张可以概括为：\textbf{开发者负责定义Agent目标、工具与边界条件，平台负责更多运行时托管职责}。这一模式的主要吸引力在于，开发者无需从零自建完整的Agent运行时。

\subsection{Harness与CMA的关系}

在深入CMA架构之前，必须首先厘清\textbf{Harness}的概念。Harness是AI Agent架构中的调度核心，Anthropic将其定义为"模型运行所依赖的指令集与护栏"。具体而言，Harness是持续调用LLM并将模型产生的工具调用路由到对应执行基础设施的控制循环。

Harness的本质是"薄调度层"（thin harness）。Anthropic的工程哲学认为，所有规划、推理和决策应推给模型，Harness本身保持"愚蠢"——只负责组装提示词、调用模型、解析输出、分发工具调用。

从工程视角看，CMA可以被理解为\textbf{对传统 Harness 模式的一种托管化封装}，但这一表述更多是本文的架构分析框架，而不应直接视为Anthropic公开文档中的正式产品术语。两者的关键差异体现在：

\begin{itemize}
    \item \textbf{自托管Harness}：开发者在自己的服务器上运行编排循环，自行维护执行环境、状态存储和凭证系统；
    \item \textbf{CMA（托管运行时）}：Anthropic托管更多运行时职责，开发者主要通过API定义Agent配置与权限边界。
\end{itemize}

\subsection{一种可供借鉴的分层理解}

为了与MHE进行结构性比较，本文采用一种\textbf{工程分析上的分层视角}来理解CMA：可将其近似拆分为会话层、编排层与执行层三个关注点。需要强调的是，这里的三层划分是本文为了分析便利而做的抽象，并不等同于Anthropic公开文档中严格承诺的官方内部实现细节。

\subsubsection{会话层（Session as control-plane state）}

根据公开文档，CMA提供了Agent与Session对象，这意味着会话状态并不只是一次性请求上下文，而是被提升为更明确的运行时对象。从工程角度，可以合理推断其具有以下价值：

\begin{itemize}
    \item \textbf{状态外部化}：会话状态不再完全绑定于单次模型调用；
    \item \textbf{上下文管理}：平台能够围绕Session组织更长时程的任务执行；
    \item \textbf{审计潜力}：会话对象为后续事件追踪和权限审计提供天然锚点。
\end{itemize}

\subsubsection{编排层（Managed orchestration）}

CMA的一个核心变化在于：开发者不再需要自建完整的Agent循环。这表明平台替开发者承担了更大比例的编排职责。由此可以抽象出以下工程含义：

\begin{itemize}
    \item \textbf{托管编排}：平台负责更多执行流程控制；
    \item \textbf{运行时封装}：开发者更多描述目标和边界，而非实现完整控制循环；
    \item \textbf{弹性潜力}：托管化设计通常更容易支持扩缩容和故障替换。
\end{itemize}

\subsubsection{执行层（Sandbox / tool execution plane）}

Anthropic公开工程文章强调了 managed agents 中“brain”与“hands”的解耦，这为理解执行层提供了启发。结合公开权限策略与凭证文档，可以合理认为：

\begin{itemize}
    \item \textbf{执行面独立化}：工具执行与控制逻辑之间存在更清晰的边界；
    \item \textbf{权限前置化}：执行能力受权限策略约束，而非完全由agent自由决定；
    \item \textbf{安全隔离增强}：凭证管理与执行面之间存在更严格的隔离设计。
\end{itemize}

因此，本文后续把CMA概括为“更偏分层、托管、解耦”的Agent基础设施实践，而不对其内部实现细节作超出公开资料的断言。

\subsection{企业级安全设计}

从公开资料能够较明确确认的是，CMA在安全方面强调权限策略与凭证保险库能力。较为稳妥的表述包括：

\begin{enumerate}
    \item \textbf{凭证保险库}：公开API文档提供了 vault credentials 相关接口；
    \item \textbf{细粒度权限策略}：公开文档包含按工具配置\texttt{always\_allow}与\texttt{always\_ask}的权限策略；
    \item \textbf{执行面受控}：托管运行时意味着执行权限并非完全由应用侧自由拼装；
    \item \textbf{更强隔离倾向}：从工程文章可见，Anthropic强调将推理控制与执行能力进行更明确分离。
\end{enumerate}

需要说明的是，诸如具体令牌注入路径、内部代理拓扑以及组织级作用域分配等实现细节，若无对应官方文档逐项佐证，则更适合作为工程推断而非事实断言。

\subsection{故障恢复与长时程执行}

从托管Agent产品形态来看，CMA显然面向长于单次API调用的任务执行场景，因此可以合理推断其必然包含某种会话恢复、错误处理与运行时替换机制。不过，就本文核对到的公开资料而言，诸如“自动检查点粒度”“\texttt{wake(sessionId)}调用语义”这类表述不宜当作已证实的公开事实。更稳妥的结论是：

\begin{itemize}
    \item \textbf{CMA支持会话化执行}，因此其运行时恢复能力应强于普通一次性请求；
    \item \textbf{托管运行时通常包含错误处理与恢复机制}，但具体实现细节未必完全公开；
    \item \textbf{对MHE的启发}并不依赖Anthropic公开全部内部机制，而在于会话化对象和托管执行这一产品方向本身。
\end{itemize}

\subsection{可观测性与审计}

公开资料显示，CMA支持面向会话和权限策略的管理能力；从工程角度看，这天然有利于形成更好的审计与观测基础。若进一步结合流式事件接口或会话事件记录，则可形成较完整的Agent运行时可观测面。但“完整事件流”“不可变审计日志”等说法，在缺乏逐项公开文档支撑时，更适合作为本文的工程化目标，而非直接归为CMA的已证实实现细节。

%==============================================================================
\section{MHE与CMA的对比分析}\label{sec:compare}
%==============================================================================

\subsection{架构维度对比}

表~\ref{tab:compare}从八个关键维度对比了MHE和CMA的架构差异。

\begin{table}[H]
\centering
\caption{MHE与CMA架构维度对比}\label{tab:compare}
\begin{tabular}{@{}p{3cm}p{5.5cm}p{5.5cm}@{}}
\toprule
\textbf{维度} & \textbf{MHE当前状态} & \textbf{CMA目标状态} \\
\midrule
状态持久化 & \texttt{GraphVersionStore}和\texttt{ComponentRegistry}均为内存实现，进程崩溃导致状态丢失 & 持久化的追加型事件日志，支持跨进程会话恢复 \\
\addlinespace
Harness状态 & \texttt{HarnessRuntime}持有活动组件对象字典，与进程强绑定 & 无状态编排器，任意实例可服务任意会话 \\
\addlinespace
沙箱隔离 & 层级抽象存在，但默认回退到进程内适配器（\texttt{InProcessAdapter}） & 惰性容器初始化，真正的进程/网络隔离 \\
\addlinespace
凭证管理 & \texttt{InMemoryIdentityBoundary}仅提供基本证明签发 & Credential Vault，作用域令牌，架构级防泄露 \\
\addlinespace
故障恢复 & 支持手动回滚到先前图版本，但无自动检查点 & 自动检查点+\texttt{wake(session\_id)}恢复 \\
\addlinespace
可观测性 & \texttt{trace\_id}仅作为关联字段，指标存储在内存中 & 完整结构化事件流，会话日志即审计源 \\
\addlinespace
水平扩展 & 单进程运行时，无法水平扩展 & 任意实例可服务任意会话，天然水平扩展 \\
\addlinespace
资源效率 & 组件长期驻留内存，沙箱持续运行 & 惰性初始化，按需分配，用完即释放 \\
\bottomrule
\end{tabular}
\end{table}

\subsection{设计理念差异}

\subsubsection{单体vs分层}

MHE当前的设计将引导、状态管理、路由、安全和可观测性等功能 bundling 在一个\texttt{HarnessRuntime}实例中。这种"单体"设计使每个会话成为需要精心照料的"宠物"（pet）：进程崩溃则会话丢失，调试困难，空闲时仍需支付全额资源成本。

CMA则通过三层解耦将Agent基础设施从"宠物"转变为可任意替换的"牲口"（cattle）——符合企业级系统对弹性和可扩展性的根本要求。

\subsubsection{功能丰富vs基础设施化}

MHE在功能层面已经非常丰富：候选图验证、热加载迁移、安全链四层防御、Pareto优化器等。然而，这些功能仍然运行在一个"研究和工程运行时"的上下文中，缺乏企业级基础设施所需的可扩展性、持久性和可观测性。

从公开定价文档看，CMA除了模型 token 成本外，还引入了按活跃会话时长计费的运行时维度。由此可以更稳妥地得出：CMA所售卖的不再只是单次推理能力，也包含持续运行的托管运行时价值。

\subsubsection{补丁式安全vs架构式安全}

MHE的安全链（沙箱验证→A/B影子测试→策略否决→自动回滚）是功能层面的安全防御。然而，其身份边界和沙箱隔离仍然存在实现上的薄弱点（如进程内回退适配器）。

CMA的安全是\textbf{架构解耦的自然结果}：凭证存储与Sandbox执行环境物理隔离，即使发生提示注入攻击，攻击者也无法通过Sandbox泄露API密钥。

%==============================================================================
\section{基于CMA设计思想的MHE改进方案}\label{sec:improvements}
%==============================================================================

基于上述对比分析，本节提出七项具体的架构改进方案。这些方案按照从基础设施到上层功能的依赖顺序排列。

\subsection{改进一：引入SessionStore——外部化运行时状态}

\subsubsection{问题定义}

当前\texttt{GraphVersionStore}、事件总线和组件状态均为内存实现：

\begin{lstlisting}[caption={当前内存中的版本存储}]
class GraphVersionStore(BaseModel):
    state: GraphState = Field(default_factory=GraphState)
    snapshots: dict[int, GraphSnapshot] = Field(default_factory=dict)
    archived: dict[int, GraphSnapshot] = Field(default_factory=dict)
    candidates: list[CandidateRecord] = Field(default_factory=list)
    retention: int = 50
\end{lstlisting}

进程崩溃后，所有图版本、候选记录和生命周期状态全部丢失。

\subsubsection{设计方案}

创建\textbf{追加型}的\texttt{SessionStore}，作为会话的单一事实源：

\begin{lstlisting}[caption={提议的SessionStore接口}]
class SessionEvent(BaseModel):
    event_id: str
    session_id: str
    timestamp: datetime
    type: EventType  # GRAPH_COMMITTED, PAYLOAD_ROUTED,
                     # SAFETY_GATE_PASSED, CHECKPOINT_SAVED, ...
    payload: dict[str, Any]

class SessionStore(Protocol):
    def append(self, session_id: str, event: SessionEvent) -> None: ...
    def get_events(self, session_id: str,
                   after: int | None = None) -> list[SessionEvent]: ...
    def slice(self, session_id: str,
              start: int, end: int) -> list[SessionEvent]: ...
    def get_checkpoint(self, session_id: str) -> int: ...
\end{lstlisting}

\texttt{SessionStore}应支持多种后端实现：
\begin{itemize}
    \item \texttt{InMemorySessionStore}：用于开发和测试；
    \item \texttt{FileSessionStore}：基于追加写文件的本地持久化；
    \item \texttt{RedisSessionStore}：基于Redis Stream的分布式实现；
    \item \texttt{SQLSessionStore}：基于关系型数据库的事务化实现。
\end{itemize}

\subsubsection{预期收益}

\begin{enumerate}
    \item \textbf{状态持久化}：进程崩溃不会丢失会话状态；
    \item \textbf{水平扩展}：多个\texttt{HarnessRuntime}实例可共享同一个\texttt{SessionStore}；
    \item \textbf{审计合规}：不可变事件日志天然满足合规要求；
    \item \textbf{调试能力}：完整的事件历史可用于事后分析和反事实诊断。
\end{enumerate}

\subsection{改进二：HarnessRuntime无状态化}

\subsubsection{问题定义}

当前\texttt{HarnessRuntime}在\texttt{boot.py}中持有活动组件实例：

\begin{lstlisting}[caption={当前有状态的HarnessRuntime}]
class HarnessRuntime:
    def __init__(self, ...):
        self.components: dict[str, HarnessComponent] = {}
        self.engine = ConnectionEngine(self.registry, ...)
        self.event_bus = EventBus()
        ...

    def boot(self) -> BootReport:
        ...
        self.components[component_id] = component
        self.engine.register_handler(record.target, handler)
        ...
\end{lstlisting}

\texttt{self.components}将组件实例与运行时进程强绑定，无法实现进程间共享或快速故障转移。

\subsubsection{设计方案}

将\texttt{HarnessRuntime}重构为\textbf{无状态编排器}：

\begin{enumerate}
    \item \textbf{移除组件实例缓存}：运行时仅持有声明和注册表条目，不缓存活动对象；
    \item \textbf{延迟激活}：组件在首次收到路由消息时才被实例化和激活；
    \item \textbf{处理器工厂模式}：将\path{register_handler(target, handler)}改为\path{register_factory(target, factory)}，在路由时动态创建处理器；
    \item \textbf{状态外部化}：组件的内部状态通过\texttt{SessionStore}持久化，而非保存在内存中。
\end{enumerate}

\begin{lstlisting}[caption={无状态HarnessRuntime示意}]
class StatelessHarnessRuntime:
    def __init__(self, session_store: SessionStore, ...):
        self.session_store = session_store
        self.registry = ComponentRegistry()
        self.engine = ConnectionEngine(self.registry, ...)
        # 不再持有组件实例

    def boot(self) -> BootReport:
        # 仅注册声明，不激活组件实例
        self.registry.register(component_id, manifest, declarations)
        ...

    def wake(self, session_id: str) -> None:
        """从SessionStore恢复会话状态。"""
        events = self.session_store.get_events(session_id)
        for event in events:
            self._replay_event(event)
\end{lstlisting}

\subsubsection{预期收益}

\begin{enumerate}
    \item \textbf{弹性}：任意\texttt{HarnessRuntime}实例可接管任意会话；
    \item \textbf{资源效率}：无活跃路由的组件不占用内存；
    \item \textbf{快速恢复}：\texttt{wake(session\_id)}可在秒级恢复会话；
    \item \textbf{水平扩展}：通过负载均衡在多个实例间分配会话。
\end{enumerate}

\subsection{改进三：惰性沙箱初始化}

\subsubsection{问题定义}

当前沙箱层级存在，但默认回退到\texttt{InProcessAdapter}：

\begin{lstlisting}[caption={当前进程内沙箱回退}]
@dataclass
class InProcessAdapter:
    tier: SandboxTier = SandboxTier.V8_WASM

    def execute(self, code, *args, **kwargs):
        try:
            output = code(*args, **kwargs)  # 直接在进程内执行
            return SandboxExecutionResult(tier=self.tier,
                                          success=True, output=output)
        except Exception as exc:
            return SandboxExecutionResult(tier=self.tier,
                                          success=False, error=str(exc))
\end{lstlisting}

这\textbf{不是}真正的安全边界，攻击者一旦突破代码限制即可访问宿主进程的全部资源。

\subsubsection{设计方案}

实现\textbf{惰性进程/容器沙箱}：

\begin{lstlisting}[caption={惰性进程沙箱适配器}]
@dataclass
class LazyProcessAdapter:
    tier: SandboxTier = SandboxTier.GVISOR
    retain: bool = False

    def _ensure_process(self):
        if self._process is None or not self._process.is_alive():
            self._process = SandboxProcess.start(tier=self.tier)
        return self._process

    def execute(self, code, *args, **kwargs):
        proc = self._ensure_process()  # 惰性初始化
        try:
            result = proc.run(code, *args, **kwargs)
            return SandboxExecutionResult(tier=self.tier,
                                          success=True, output=result)
        finally:
            if not self.retain:
                proc.destroy()
                self._process = None
\end{lstlisting}

同时引入\textbf{沙箱资源配额}管理：
\begin{itemize}
    \item CPU时间限制；
    \item 内存使用上限；
    \item 网络访问白名单；
    \item 文件系统只读/读写权限控制。
\end{itemize}

\subsubsection{预期收益}

\begin{enumerate}
    \item \textbf{真实隔离}：代码执行在独立进程中，即使崩溃也不影响Harness；
    \item \textbf{资源效率}：仅在需要时分配沙箱资源；
    \item \textbf{弹性}：单个沙箱崩溃可快速重建，不影响整体服务；
    \item \textbf{合规性}：满足企业级安全审计对执行隔离的要求。
\end{enumerate}

\subsection{改进四：凭证保险库（Credential Vault）}

\subsubsection{问题定义}

当前\texttt{InMemoryIdentityBoundary}仅提供基本的证明签发功能：

\begin{lstlisting}[caption={当前最小化身份边界}]
class InMemoryIdentityBoundary:
    def issue_attestation(self, subject: str) -> IdentityAttestation: ...
    def sanitize_payload(self, payload: Any) -> Any: ...
    def get_credentials(self, attestation_id: str) -> Any: ...
\end{lstlisting}

凭证仍可能通过\texttt{sanitize\_payload}泄漏到普通载荷路径，且缺乏细粒度的权限策略。

\subsubsection{设计方案}

引入\textbf{Credential Vault}，实现架构级凭证隔离：

\begin{lstlisting}[caption={凭证保险库设计}]
class CredentialVault:
    """凭证保险库：凭证仅写不可读，不进入Prompt、日志或Sandbox。"""

    def store(self, key: str, credential: str,
              scope: str | None = None) -> CredentialRef:
        """存储凭证，返回不可逆的引用句柄。"""
        ...

    def inject(self, ref: CredentialRef,
               target_sandbox: SandboxAdapter) -> None:
        """将凭证注入目标沙箱，Harness本身不接触明文。"""
        ...

    def get_public_view(self, ref: CredentialRef) -> SanitizedCredential:
        """返回去敏化的凭证视图，供日志和审计使用。"""
        ...

    def set_policy(self, tool_name: str,
                   policy: Literal["always_allow", "always_ask"]) -> None:
        """按工具配置权限策略。"""
        ...

    def revoke(self, ref: CredentialRef) -> None:
        """实时撤销凭证。"""
        ...
\end{lstlisting}

关键安全属性：
\begin{enumerate}
    \item \textbf{写不可读}：凭证一旦存入，无法通过API读取明文；
    \item \textbf{作用域隔离}：不同会话/工具的凭证相互隔离；
    \item \textbf{外部注入}：凭证在沙箱初始化时通过安全通道注入，不出现在任何中间层；
    \item \textbf{审计追踪}：所有凭证操作（存储、注入、撤销）记录到\texttt{SessionStore}。
\end{enumerate}

\subsubsection{预期收益}

\begin{enumerate}
    \item \textbf{防泄露}：即使发生提示注入攻击，攻击者也无法获取原始凭证；
    \item \textbf{合规}：满足企业对凭证管理的审计要求；
    \item \textbf{权限治理}：支持按工具、按会话的细粒度权限控制。
\end{enumerate}

\subsection{改进五：自动检查点与故障恢复}

\subsubsection{问题定义}

MHE当前支持手动回滚到先前图版本，但缺乏自动化的检查点和故障恢复机制。当\texttt{HarnessRuntime}进程崩溃时，所有进行中的会话状态丢失，必须从头重建。

\subsubsection{设计方案}

在\texttt{SessionStore}之上实现\textbf{自动检查点}机制：

\begin{lstlisting}[caption={自动检查点管理器}]
class CheckpointManager:
    def __init__(self, session_store: SessionStore,
                 interval: int = 10):
        self.store = session_store
        self.interval = interval  # 每N个事件打一个检查点

    def maybe_checkpoint(self, session_id: str) -> None:
        events = self.store.get_events(session_id)
        if len(events) % self.interval == 0:
            checkpoint = SessionCheckpoint(
                session_id=session_id,
                event_count=len(events),
                graph_version=self._get_active_graph(events),
                component_states=self._capture_states(events),
            )
            self.store.save_checkpoint(session_id, checkpoint)

    def recover(self, session_id: str) -> HarnessState:
        checkpoint = self.store.get_latest_checkpoint(session_id)
        if checkpoint is None:
            raise ValueError(f"No checkpoint found for {session_id}")

        # 从检查点恢复基础状态
        state = HarnessState.from_checkpoint(checkpoint)

        # 重放检查点之后的事件
        events = self.store.get_events(session_id,
                                       after=checkpoint.event_count)
        for event in events:
            state.apply_event(event)

        return state
\end{lstlisting}

\textbf{关键检查点触发时机}：
\begin{itemize}
    \item 图提交成功后；
    \item 安全链全部通过后；
    \item 关键工具执行完成后；
    \item 显式调用\texttt{checkpoint()}时。
\end{itemize}

\subsubsection{预期收益}

\begin{enumerate}
    \item \textbf{自动恢复}：进程崩溃后，新实例可在秒级恢复会话；
    \item \textbf{长时程任务}：支持运行数小时甚至数天的Agent任务；
    \item \textbf{可靠性}：从"尽人事"变为"制度化保障"。
\end{enumerate}

\subsection{改进六：事件驱动的可观测性}

\subsubsection{问题定义}

当前\texttt{trace\_id}主要作为关联字段，指标存储在内存中，缺乏结构化的事件流：

当前trace\_id仅用于审计事件关联，不改变运行时行为。

\begin{lstlisting}[caption={}]
# trace_id 仅用于审计事件关联，不改变运行时行为
ObservabilityComponent.record_event(..., trace_id=trace_id)
\end{lstlisting}

\subsubsection{设计方案}

建立基于\texttt{SessionStore}的\textbf{统一事件流}：

\begin{lstlisting}[caption={结构化事件类型}]
class EventType(str, Enum):
    # 图生命周期事件
    GRAPH_CANDIDATE_CREATED = "graph.candidate.created"
    GRAPH_VALIDATED = "graph.validated"
    GRAPH_COMMITTED = "graph.committed"
    GRAPH_ROLLED_BACK = "graph.rolled_back"

    # 路由事件
    PAYLOAD_RECEIVED = "payload.received"
    PAYLOAD_ROUTED = "payload.routed"
    HANDLER_EXECUTED = "handler.executed"
    HANDLER_FAILED = "handler.failed"

    # 安全事件
    SAFETY_GATE_PASSED = "safety.gate.passed"
    SAFETY_GATE_REJECTED = "safety.gate.rejected"
    SANDBOX_EXECUTED = "sandbox.executed"

    # 系统事件
    COMPONENT_ACTIVATED = "component.activated"
    COMPONENT_SUSPENDED = "component.suspended"
    CHECKPOINT_SAVED = "checkpoint.saved"
    SESSION_WAKED = "session.waked"
\end{lstlisting}

提供\textbf{SSE式的事件流接口}：

\begin{lstlisting}[caption={事件流接口}]
class EventStream:
    def subscribe(self, session_id: str,
                  callback: Callable[[SessionEvent], None]) -> str:
        """订阅指定会话的事件流。"""
        ...

    def unsubscribe(self, subscription_id: str) -> None:
        ...

    def query(self, session_id: str,
              event_types: list[EventType] | None = None,
              start_time: datetime | None = None,
              end_time: datetime | None = None) -> list[SessionEvent]:
        """结构化查询历史事件。"""
        ...
\end{lstlisting}

\subsubsection{预期收益}

\begin{enumerate}
    \item \textbf{实时可观测}：通过SSE推送实时了解Agent思考过程；
    \item \textbf{审计合规}：完整的事件历史满足合规要求；
    \item \textbf{调试效率}：结构化查询替代日志检索；
    \item \textbf{反事实分析}：基于事件历史进行what-if分析。
\end{enumerate}

\subsection{改进七：BrainProvider抽象——解耦模型与Harness}

\subsubsection{问题定义}

当前\texttt{ComponentRuntime.llm}字段类型为\texttt{Any}，Harness与特定模型后端隐式耦合：

\begin{lstlisting}[caption={当前模糊的LLM接口}]
@dataclass
class ComponentRuntime:
    llm: Any | None = None  # 类型不透明
    ...
\end{lstlisting}

\subsubsection{设计方案}

引入显式的\texttt{BrainProvider}抽象：

\begin{lstlisting}[caption={BrainProvider抽象}]
class Plan(BaseModel):
    tool_calls: list[ToolCall]
    reasoning: str | None = None

class BrainProvider(Protocol):
    """模型提供者抽象：Harness仅传递上下文，不干预决策。"""

    async def plan(self,
                   context: list[SessionEvent],
                   tools: list[Tool],
                   model: str | None = None) -> Plan:
        """基于上下文和可用工具生成执行计划。"""
        ...

    def get_default_model(self) -> str:
        ...

    def estimate_cost(self, context: list[SessionEvent]) -> CostEstimate:
        ...

class AdvisorLikeStrategy:
    """顾问策略：廉价模型咨询昂贵模型，降低长时程任务成本。"""

    def __init__(self,
                 executor_brain: BrainProvider,  # 如 Haiku
                 advisor_brain: BrainProvider):   # 如 Opus
        self.executor = executor_brain
        self.advisor = advisor_brain

    async def plan(self, context, tools):
        # 先由廉价模型尝试
        plan = await self.executor.plan(context, tools)
        # 在复杂决策点咨询昂贵模型
        if self._needs_advice(plan):
            plan = await self.advisor.plan(context, tools)
        return plan
\end{lstlisting}

\subsubsection{预期收益}

\begin{enumerate}
    \item \textbf{模型无关}：Harness不依赖特定模型后端；
    \item \textbf{成本优化}：可进一步在该抽象之上叠加分层模型选择或顾问式推理策略；
    \item \textbf{可测试性}：BrainProvider可被Mock替换，便于单元测试。
\end{enumerate}

%==============================================================================
\section{实施路线图}\label{sec:roadmap}
%==============================================================================

\subsection{阶段划分}

基于依赖关系和实现复杂度，将七项改进划分为三个阶段实施，如表~\ref{tab:roadmap}所示。

\begin{table}[H]
\centering
\caption{分阶段实施路线图}\label{tab:roadmap}
\begin{tabular}{@{}clp{8cm}@{}}
\toprule
\textbf{阶段} & \textbf{时间} & \textbf{工作内容} \\
\midrule
\multirow{3}{*}{P1} & \multirow{3}{*}{1-2个月} & 1. 设计并实现\texttt{SessionStore}抽象及内存/文件后端 \\
 & & 2. 定义\texttt{SessionEvent}事件模式 \\
 & & 3. 将\texttt{HarnessRuntime}重构为无状态，实现\texttt{wake()}恢复 \\
\midrule
\multirow{3}{*}{P2} & \multirow{3}{*}{2-3个月} & 4. 实现惰性进程沙箱适配器（替代\texttt{InProcessAdapter}） \\
 & & 5. 实现\texttt{CredentialVault}及按工具权限策略 \\
 & & 6. 实现\texttt{CheckpointManager}自动检查点机制 \\
\midrule
\multirow{2}{*}{P3} & \multirow{2}{*}{1-2个月} & 7. 实现\texttt{EventStream} SSE式事件流 \\
 & & 8. 实现\texttt{BrainProvider}抽象及顾问式推理策略 \\
\bottomrule
\end{tabular}
\end{table}

\subsection{关键文件修改清单}

\begin{table}[H]
\centering
\caption{关键文件修改清单}\label{tab:files}
\begin{tabular}{@{}lll@{}}
\toprule
\textbf{文件路径} & \textbf{修改类型} & \textbf{关联改进} \\
\midrule
\texttt{metaharness/core/session\_store.py} & 新增 & P1-改进一 \\
\texttt{metaharness/observability/events.py} & 新增 & P1-改进一/P3-改进六 \\
\texttt{metaharness/core/boot.py} & 重构 & P1-改进二/P2-改进五 \\
\texttt{metaharness/core/connection\_engine.py} & 修改 & P1-改进一 \\
\texttt{metaharness/safety/sandbox\_tiers.py} & 重构 & P2-改进三 \\
\texttt{metaharness/safety/sandbox\_process.py} & 新增 & P2-改进三 \\
\texttt{metaharness/identity/vault.py} & 新增 & P2-改进四 \\
\texttt{metaharness/core/checkpoint.py} & 新增 & P2-改进五 \\
\texttt{metaharness/observability/stream.py} & 新增 & P3-改进六 \\
\texttt{metaharness/providers/brain.py} & 新增 & P3-改进七 \\
\bottomrule
\end{tabular}
\end{table}

\subsection{风险与缓解措施}

\begin{table}[H]
\centering
\caption{风险与缓解措施}\label{tab:risks}
\begin{tabular}{@{}p{3.5cm}p{5cm}p{5cm}@{}}
\toprule
\textbf{风险} & \textbf{影响} & \textbf{缓解措施} \\
\midrule
状态持久化引入延迟 & 事件追加写入可能降低路由吞吐量 & 采用异步追加+内存缓冲的混合策略；支持批量写入 \\
无状态化破坏向后兼容 & 现有组件可能依赖运行时中的实例缓存 & 提供兼容层\texttt{StatefulHarnessRuntime}，标记为废弃 \\
惰性沙箱启动开销 & 首次工具调用时延增加 & 预温热池（warm pool）机制；对高频工具保持少量常驻沙箱 \\
凭证vault增加复杂度 & 开发者学习曲线变陡 & 提供默认的内存vault用于开发；生产环境再切换到加密vault \\
\bottomrule
\end{tabular}
\end{table}

%==============================================================================
\section{结论与展望}\label{sec:conclusion}
%==============================================================================

\subsection{主要结论}

本文通过系统分析Anthropic CMA的架构设计原则，识别出MHE框架在状态持久化、Harness无状态化、沙箱隔离、凭证管理和故障恢复等五个关键维度的差距，并提出了七项具体的架构改进方案。主要结论如下：

\begin{enumerate}
    \item \textbf{基础设施价值来自解耦，而非功能}。MHE已经具备丰富的功能（候选图验证、热加载、安全链、优化器），但这些功能捆绑在单体进程中。通过引入SessionStore将状态外部化，可以使Harness无状态化，从而自然获得水平扩展、故障恢复和审计合规等能力。

    \item \textbf{安全应当是架构的自然结果}。当前MHE的安全链是功能层面的防御，而凭证隔离和沙箱边界存在实现薄弱点。通过Credential Vault和惰性进程沙箱，可以将安全从"补丁"提升为"架构属性"。

    \item \textbf{事件流是统一抽象}。SessionStore中的追加型事件日志可以同时支撑状态恢复、可观测性、审计追踪和反事实分析，避免为每个目的建立独立的数据管道。

    \item \textbf{模型应与Harness解耦}。通过BrainProvider抽象，MHE可以支持多模型后端和顾问式推理策略，避免对单一模型族的过度依赖。
\end{enumerate}

\subsection{未来研究方向}

\begin{enumerate}
    \item \textbf{多租户支持}：在SessionStore和CredentialVault基础上，实现命名空间级别的资源隔离和配额管理；
    \item \textbf{分布式图执行}：将ConnectionEngine扩展为支持跨节点路由的分布式版本；
    \item \textbf{自适应沙箱层级}：根据工具风险画像动态选择沙箱隔离强度（V8/WASM vs gVisor vs Firecracker）；
    \item \textbf{智能检查点策略}：基于事件重要性而非固定间隔决定检查点时机，减少存储开销；
    \item \textbf{与Aeloon Plugin SDK的深度集成}：将MHE作为Aeloon插件生态的基础设施层，利用Plugin SDK的声明式模型统一组件生命周期管理。
\end{enumerate}

\subsection{最终评述}

CMA的发布证明Anthropic已清醒地认识到：在模型能力趋同的时代，真正的护城河不是权重，而是让模型可靠、安全、规模化运行的基础设施。对于MHE框架而言，这意味着从"研究和工程运行时"向"生产级控制平面"的演进不仅是技术升级，更是范式转变。本文提出的改进方案为这一转变提供了明确的架构蓝图和可操作的实施路径。

%==============================================================================
% References
%==============================================================================
\newpage
\section*{参考文献}
\addcontentsline{toc}{section}{参考文献}

\begin{enumerate}[label={[\arabic*]}]
    \item Lee J, et al. Meta-Harness: Optimizing LLM Harnesses through File-System-Based Experience. \textit{Stanford/MIT}, 2026.
    \item Hu S, et al. Automated Design of Agentic Systems. \textit{UBC/Vector Institute}, 2025.
    \item Yin H, et al. Gödel Agent: A Self-Referential Self-Improvement Framework. 2024.
    \item Zhang T, et al. Darwin Gödel Machine: Open-Ended Self-Improvement through Evolution. 2025.
    \item Anthropic. Claude Managed Agents Overview / Quickstart / Onboarding / Permission Policies / Pricing / Vault Credentials. \textit{Anthropic Platform Documentation}, 2026.
    \item Anthropic. Scaling Managed Agents: Decoupling the Brain from the Hands. \textit{Anthropic Engineering Blog}, 2026.
    \item Basu A, et al. Rigorous Component-Based System Design Using the BIP Framework. \textit{IEEE Software}, 2011.
    \item Ling S. DynaComm: Dynamic Component-Based Reconfiguration. 2007.
    \item MHE Technical Manual. \texttt{MHE/docs/TECHNICAL\_MANUAL.md}, 2026.
    \item MHE User Guide. \texttt{MHE/docs/USER\_GUIDE.md}, 2026.
    \item Intra-CMA Analysis. \texttt{MHE/docs/Intra-CMA.md}, 2026.
\end{enumerate}

\end{document}
